% ============================================================
% CHAPTER 1: INTRODUCTION (matches Vazirani Ch. 1)
% ============================================================
\setcounter{chapter}{0}
\chapter{Introduction}
\label{ch:intro}

\section{The Approximation Paradigm}

\begin{intuitionbox}
\textbf{The Core Problem:} Most natural optimization problems are NP-hard. Assuming P $\neq$ NP, we cannot find exact solutions efficiently. But we \emph{can} find near-optimal solutions with \emph{provable guarantees}.

\textbf{The Key Insight:} To prove a guarantee like ``within 2$\times$ optimal,'' we don't need to \emph{know} the optimal value. We just need a \emph{lower bound} on OPT that we can compute efficiently. If our algorithm's solution is within a factor of that lower bound, we have our guarantee.
\end{intuitionbox}

\begin{definition}[$\alpha$-Approximation Algorithm]
For a minimization problem, an algorithm $A$ is an $\alpha$-approximation ($\alpha \geq 1$) if for every instance $I$:
\[
A(I) \leq \alpha \cdot \OPT(I)
\]
where $\OPT(I)$ is the optimal value. For maximization: $A(I) \geq \frac{1}{\alpha}\OPT(I)$.
\end{definition}

\subsection{Vertex Cover: A Gentle Beginning}

\begin{definition}[Vertex Cover]
Given an undirected graph $G = (V, E)$, a \textbf{vertex cover} is a subset $C \subseteq V$ such that every edge has at least one endpoint in $C$. The \textbf{cardinality vertex cover problem} asks for a cover of minimum size.
\end{definition}

\subsection{Lower Bound via Maximal Matching}

\begin{definition}[Matching]
A \textbf{matching} $M \subseteq E$ is a set of edges with no shared endpoints. It is \textbf{maximal} if no edge can be added without violating the matching property.
\end{definition}

\begin{lemma}
For any graph $G$, the size of a maximal matching $M$ is a lower bound on the minimum vertex cover: $|M| \leq \OPT$.
\end{lemma}
\begin{proof}
Each edge in $M$ needs a distinct vertex in any vertex cover (edges don't share endpoints). So $|C| \geq |M|$ for any cover $C$.
\end{proof}

\begin{intuitionbox}
\textbf{Intuition:} A maximal matching is easy to find greedily. It pairs up vertices that ``must be covered'' in some sense. Each pair contributes 2 vertices to our cover but at most 1 to the optimal cover. This 2:1 ratio is the source of the factor-2 guarantee.
\end{intuitionbox}

\subsection{The 2-Approximation Algorithm}

\begin{algorithmdef}[Vertex Cover via Maximal Matching (Factor 2)]
\begin{enumerate}
    \item Find a maximal matching $M$ in $G$ (greedily).
    \item Output $C = \{u, v \mid (u,v) \in M\}$ — all matched vertices.
\end{enumerate}
\end{algorithmdef}

\begin{theorem}
Algorithm 1.1 is a factor-2 approximation for Vertex Cover.
\end{theorem}
\begin{proof}
$C$ is a vertex cover because $M$ is maximal (any uncovered edge could be added to $M$). Since $|M| \leq \OPT$ and $|C| = 2|M|$, we have $|C| \leq 2\OPT$.
\end{proof}

\subsection{Why Approximation? Why Not Exact?}

\begin{itemize}
    \item \textbf{Exact Vertex Cover is NP-complete} (Karp, 1972). No polynomial-time exact algorithm exists unless P = NP.
    \item \textbf{Fixed-parameter algorithms} exist ($O(1.2738^k + kn)$) but exponential in $k$.
    \item \textbf{Integer Linear Programming} solves it exactly but worst-case exponential.
    \item \textbf{Approximation gives polynomial time with a guarantee} — often the best practical choice.
\end{itemize}

\subsection{Other Approaches and Their Limits}

\begin{table}[H]
\centering
\caption{Approaches to Vertex Cover}
\begin{tabular}{lcc}
\toprule
Method & Time & Guarantee \\
\midrule
Brute force / ILP & Exponential & Exact \\
Fixed-parameter & $O(1.2738^k)$ & Exact, but exponential in $k$ \\
Maximal matching (this chapter) & $O(m)$ & Factor 2 \\
LP rounding (Chapter 14) & Poly & Factor 2 \\
Semidefinite programming & Poly & Factor $2-\epsilon$ (open) \\
\bottomrule
\end{tabular}
\end{table}

\textbf{Why not better than 2?} Under the Unique Games Conjecture, no polynomial-time algorithm can achieve $(2-\epsilon)$-approximation (Khot-Regev, 2008). The factor 2 is essentially tight.

\subsection{Python Implementation}

\begin{lstlisting}[style=python, caption=Vertex Cover 2-Approximation]
def maximal_matching(graph):
    """Greedy maximal matching: pick edges, remove endpoints."""
    n = len(graph)
    matched = [False] * n
    matching = []
    for u in range(n):
        if not matched[u]:
            for v in graph[u]:
                if not matched[v]:
                    matching.append((u, v))
                    matched[u] = matched[v] = True
                    break
    return matching

def vertex_cover_approx_2(graph):
    """Factor-2 approximation for Vertex Cover."""
    matching = maximal_matching(graph)
    cover = set()
    for u, v in matching:
        cover.add(u)
        cover.add(v)
    return cover
\end{lstlisting}

\section{Applications}
\begin{itemize}
    \item \textbf{Network monitoring:} Place sensors on vertices to monitor all links
    \item \textbf{Compiler register allocation:} Variables as vertices, conflicts as edges
    \item \textbf{Bioinformatics:} Protein interaction networks — find minimal interacting set
    \item \textbf{Scheduling:} Tasks as vertices, conflicts as edges
\end{itemize}

\subsection*{Real Example: Register Allocation in a Compiler}
\textbf{Scenario:} A compiler must assign $n=50$ variables to $k=16$ CPU registers. Variables used simultaneously cannot share a register.
\textbf{Model:} Conflict graph $G=(V,E)$ where $V$ = variables, $(u,v) \in E$ if $u,v$ are live simultaneously. Minimum vertex cover finds the minimum set of variables to \emph{spill to memory} so the remaining $k$ can fit in registers.
\textbf{Why 2-approx matters:} If OPT spills 10 variables, our algorithm spills $\leq 20$. In practice, maximal matching often finds near-optimal solutions. The guarantee means we never spill more than twice the minimum — critical for register-starved architectures (embedded systems).
\textbf{Alternative approach:} ILP (exact) takes exponential time; greedy degree-based gives $O(\log n)$.

\subsection{Tight Example}

\begin{example}[Complete Bipartite $K_{n,n}$]
The algorithm picks all $2n$ vertices. Optimal picks one side ($n$ vertices). Ratio $= 2$.
\end{example}

\subsection{Exercises}
\begin{exercise}
Show that the greedy algorithm picking highest-degree vertices gives $O(\log n)$ approximation.
\end{exercise}
\begin{exercise}
Design a 2-approximation for \emph{minimum cardinality maximal matching}.
\end{exercise}

